\section{附录}

\subsection{当前实现位置}

当前 MGR-SID v1 的实现位于仓库中的独立实验分支，不会改动默认 MiniOneRec tokenizer 主链。主要入口为：
\begin{itemize}[leftmargin=1.4em]
  \item \texttt{scripts/experiment\_mgr\_sid\_v1\_train.py}
  \item \texttt{scripts/experiment\_mgr\_sid\_v1\_generate.py}
\end{itemize}

核心实现位于：
\begin{itemize}[leftmargin=1.4em]
  \item \texttt{src/onerec/experiments/mgr\_sid/train\_v1.py}
  \item \texttt{src/onerec/experiments/mgr\_sid/transplanted\_graph\_bank.py}
  \item \texttt{src/onerec/experiments/mgr\_sid/paper\_transplants.py}
\end{itemize}

\subsection{v1 中的 graph bank 设置}

\begin{table}[h]
  \centering
  \caption{当前 v1 实验中实际使用的 graph-bank 参数。}
  \label{tab:appendix-graph}
  \begin{tabular}{ll}
    \toprule
    组件 & 设定 \\
    \midrule
    History window & 10 \\
    Coarse minimum edge weight & 2.0 \\
    Local minimum edge weight & 1.0 \\
    Semantic anchor top-$k$ & 32 \\
    Semantic anchor mixing weight & 0.35 \\
    Spectral rank & 48 \\
    Middle-band eigen ratio & [0.25, 0.65] \\
    Temporal mixing ratio (alt. middle view) & 0.35 \\
    Graph top-$k$ per row in training & 32 \\
    \bottomrule
  \end{tabular}
\end{table}

\subsection{Tokenizer 训练超参数}

\begin{table}[h]
  \centering
  \caption{用于 \Cref{tab:main} 公平比较的 upstream-aligned tokenizer 训练制度。}
  \label{tab:appendix-train}
  \begin{tabular}{ll}
    \toprule
    超参数 & 数值 \\
    \midrule
    量化层数 & 3 \\
    Codebook sizes & [256, 256, 256] \\
    Quantized embedding dimension & 32 \\
    Encoder/decoder widths & [2048, 1024, 512, 256, 128, 64] \\
    Optimizer & AdamW \\
    Learning rate & $10^{-3}$ \\
    Weight decay & 0 \\
    Scheduler & constant with warmup \\
    Warmup epochs & 50 \\
    Epochs & 10000 \\
    Batch size & 20480 \\
    Quantization loss weight & 1.0 \\
    $\beta$ & 0.25 \\
    Graph loss weights & $(0.05, 0.15, 0.05)$ \\
    \bottomrule
  \end{tabular}
\end{table}

\subsection{局部歧义分析的解释}

本文中的 local ambiguity analysis 比较的是 final baseline index 与 final hierarchy index，而不是直接比较推荐输出。这是一个 tokenizer-native 的分析，因此应被理解为对 tokenizer 结构本身的证据：
\begin{itemize}[leftmargin=1.4em]
  \item hierarchy-aware tokenizer 产生了更少拥挤的 same-$l_2$ bucket；
  \item 这一变化在 catalog 级和 test-weighted 两种口径下都成立；
  \item 它与本文的核心动机，即 local leaf ambiguity，直接对应。
\end{itemize}

\subsection{当前稿件的范围}

这份稿件应被视为 tokenizer-stage draft，而不是完整 end-to-end paper。接下来最直接的里程碑包括：
\begin{enumerate}[leftmargin=1.4em]
  \item 对 final tokenizer comparison 补做 repeated seed 验证；
  \item 在 Office 上补齐 upstream-aligned tokenizer 实验；
  \item 将新的 hierarchy-aware SID 推进到 SFT 再到 RL；
  \item 等下游指标具备后，再与更多已发表 tokenizer 方法做主表比较。
\end{enumerate}
